% Part A (tech core) of the AI-for-X seminars: what AI, ML, deep learning and generative AI are, for a technical
% audience. 9 frames. Expanded 2026-11-07 at the author's request ("all the technical slides should be amply there"):
% definitions, hierarchy, traditional vs ML, types of ML, deep learning, next-token prediction, transformers,
% generative AI across media, foundation models. Some frames adapted from ai_intro_tech_leaders_ai101 and
% ai_intro_tech_leaders_llmagents. The non-tech core stays at 16 frames; the tech core is deliberately larger.
% Part A is shared: edit here, then recompile every deck that includes it. Trim by commenting frames, not deleting.

%%%%%%%%%%%%%%%%%%%%%%%%%%%%%%%%%%%%%%%%%%%%%%%%%%%%%%%%%%
\begin{frame}[fragile]\frametitle{What Is AI? A Working Definition}
\textbf{Artificial Intelligence (AI):}
\begin{itemize}
\item Machines performing tasks that \textit{normally require human intelligence}: perceiving, recognising patterns, reasoning, deciding, generating content.
\item \textbf{My definition:} ``If a computer program starts doing intelligent tasks, perceiving, reasoning, deciding, creating, that's AI.''
\end{itemize}
\textbf{How it differs from traditional software:}
\begin{itemize}
\item Traditional software: \textit{explicit rules} written by a programmer produce the output.
\item AI/ML: \textit{data + a learning algorithm}; the rules emerge from the data.
\end{itemize}
\end{frame}

%%%%%%%%%%%%%%%%%%%%%%%%%%%%%%%%%%%%%%%%%%%%%%%%%%%%%%%%%%
\begin{frame}[fragile]\frametitle{The Hierarchy: AI, ML, Deep Learning, Generative AI}
\begin{columns}
    \begin{column}[T]{0.58\linewidth}
      \begin{itemize}
      \item \textbf{AI (broadest):} any system doing something that looks intelligent, hand-written rules included.
      \item \textbf{ML (subset of AI):} behaviour learned from data instead of programmed.
      \item \textbf{Deep learning (subset of ML):} neural networks with many layers; powers vision, speech and language.
      \item \textbf{Generative AI (subset of DL):} models that create new text, code, images or audio.
      \item \textbf{LLMs:} generative models for text, trained on very large corpora.
      \item \textbf{Agents:} LLMs given tools and memory so they can take actions toward a goal.
      \end{itemize}
    \end{column}
    \begin{column}[T]{0.4\linewidth}
      \begin{center}
      \includegraphics[width=\linewidth,keepaspectratio]{ai_hierarchy_updated}
      \end{center}
    \end{column}
\end{columns}
\end{frame}

%%%%%%%%%%%%%%%%%%%%%%%%%%%%%%%%%%%%%%%%%%%%%%%%%%%%%%%%%%
\begin{frame}[fragile]\frametitle{Traditional Programming vs.\ Machine Learning}
\begin{center}
\includegraphics[width=0.45\linewidth,keepaspectratio]{tradml}
\end{center}
\begin{itemize}
\item ML inverts the loop: you supply data and the desired outputs, and the training algorithm finds the program, a model with many numeric parameters.
\item Training adjusts those parameters to reduce error. Inference runs the fixed model on new inputs.
\item The real test of a model is data it has not seen during training.
\end{itemize}
\end{frame}

%%%%%%%%%%%%%%%%%%%%%%%%%%%%%%%%%%%%%%%%%%%%%%%%%%%%%%%%%%
\begin{frame}[fragile]\frametitle{Types of Machine Learning}
\begin{itemize}
\item \textbf{Supervised learning:} learn from \textit{labelled} examples (input plus correct answer).
  \begin{itemize}
  \item Example: email to spam or not spam; scan image to tumour or benign.
  \end{itemize}
\item \textbf{Unsupervised learning:} find structure in \textit{unlabelled} data.
  \begin{itemize}
  \item Example: cluster customers into segments with no predefined groups.
  \end{itemize}
\item \textbf{Reinforcement learning:} an agent learns from \textit{rewards and penalties} for its actions.
  \begin{itemize}
  \item Example: AlphaGo learned Go partly by playing millions of games against itself; AlphaZero later learned chess the same way.
  \end{itemize}
\item \textbf{Self-supervised learning:} the labels come from the data itself, such as predicting a hidden or next word. This is how LLMs are pretrained.
\end{itemize}
\end{frame}

%%%%%%%%%%%%%%%%%%%%%%%%%%%%%%%%%%%%%%%%%%%%%%%%%%%%%%%%%%
\begin{frame}[fragile]\frametitle{Deep Learning: Neural Networks}
\begin{columns}
    \begin{column}[T]{0.52\linewidth}
      \begin{itemize}
      \item A neural network passes numbers through layers of simple units. Each connection has a weight.
      \item Training adjusts all the weights (backpropagation with gradient descent) so the output gets closer to the target.
      \item With many layers, the network learns its own features: edges, then shapes, then objects; or characters, then words, then meaning.
      \item It took off around 2012 when three things met: large datasets, GPUs, and better training methods.
      \end{itemize}
    \end{column}
    \begin{column}[T]{0.45\linewidth}
      \begin{center}
      \includegraphics[width=\linewidth,keepaspectratio]{dl25}
      \end{center}
    \end{column}
\end{columns}
\end{frame}

%%%%%%%%%%%%%%%%%%%%%%%%%%%%%%%%%%%%%%%%%%%%%%%%%%%%%%%%%%
\begin{frame}[fragile]\frametitle{Language Models: Predict the Next Token}
\begin{itemize}
\item A \textbf{language model} gives a probability for each possible next token, given the text so far: $P(w_{n+1} \mid w_1 \ldots w_n)$.
\item Generation repeats one step: pick a token from that distribution, append it, predict again.
\item An \textbf{LLM} is such a model with billions of parameters, trained on a very large amount of text.
\item \textbf{Pretraining} on raw text gives fluency and broad knowledge. Later stages (instruction tuning, preference tuning) teach it to follow requests and answer helpfully.
\end{itemize}
\end{frame}

%%%%%%%%%%%%%%%%%%%%%%%%%%%%%%%%%%%%%%%%%%%%%%%%%%%%%%%%%%
\begin{frame}[fragile]\frametitle{How Transformers Work: Attention}
\begin{columns}
    \begin{column}[T]{0.6\linewidth}
      \begin{itemize}
      \item Older recurrent networks (RNNs) read text one word at a time and struggled to keep early context.
      \item A \textbf{Transformer} (2017, ``Attention Is All You Need'') processes the whole sequence at once and computes how much each token should attend to every other token.
      \item $\text{Attention}(Q, K, V) = \text{softmax}\!\left(\frac{QK^T}{\sqrt{d_k}}\right)V$
      \item In plain words: score every word against every other word, then mix their meanings by those scores. ``Bank'' in ``river bank'' attends more to ``river'' than to ``money.''
      \item Almost every current LLM is a stack of Transformer layers.
      \end{itemize}
    \end{column}
    \begin{column}[T]{0.37\linewidth}
      \begin{center}
      \includegraphics[width=\linewidth,keepaspectratio]{transformer}
      \end{center}
    \end{column}
\end{columns}
\end{frame}

%%%%%%%%%%%%%%%%%%%%%%%%%%%%%%%%%%%%%%%%%%%%%%%%%%%%%%%%%%
\begin{frame}[fragile]\frametitle{Generative AI: Text, Code, Images, Audio}
\begin{itemize}
\item \textbf{Text and code:} LLMs generate token by token. Code is just another kind of text to them.
\item \textbf{Images:} most image generators are \textbf{diffusion models}. They start from random noise and remove it step by step, steered by the text prompt.
\item \textbf{Audio and speech:} models generate or transcribe sound in the same learn-from-examples way.
\item \textbf{Multimodal models:} one model takes in and produces several kinds of data, for example reading an image and answering in text.
\item The common idea: learn the patterns of a very large set of examples, then produce a new example that fits them.
\end{itemize}
\end{frame}

%%%%%%%%%%%%%%%%%%%%%%%%%%%%%%%%%%%%%%%%%%%%%%%%%%%%%%%%%%
\begin{frame}[fragile]\frametitle{Foundation Models and How to Adapt Them}
\begin{itemize}
\item A \textbf{foundation model} is trained once on broad data, then adapted to many tasks without building a new model each time.
\item Three ways to adapt it, from cheapest to costliest: write a better prompt, add retrieved documents to the prompt (RAG), or fine-tune the weights.
\item Chat assistants and coding assistants are products built on these models. The model is the engine; the product adds tools, memory and guardrails around it.
\item For most organisations the work is integrating existing models well, not training new ones.
\end{itemize}
\end{frame}
